% chapters/chapter7_conclusion.tex
\chapter{Conclusion}
\label{ch:conclusion}

% ============================================================
% CHAPTER OVERVIEW
% ============================================================
% This chapter wraps up the thesis:
% - Summary of findings
% - Answers to research questions
% - Contributions
% - Final remarks
%
% TARGET: 4-6 pages total
%
% STYLE NOTES:
% - Be concise and direct
% - Don't introduce new information
% - Reference specific findings from Results/Discussion


% ============================================================
% 7.1 SUMMARY OF FINDINGS
% ============================================================
\section{Summary of Findings}
\label{sec:concl_summary}

% TARGET: 1.5-2 pages
%
% CONTENT:
%
% Brief recap of the problem:
% - Apple Catcher achieved 78% cross-subject accuracy
% - Pre-cue window dramatically outperformed post-cue
% - Unclear if this reflected legitimate neural signal or artifacts
%
% What we found:
%
% Regarding the neural content of the data:
% - [Summary of ERDS findings - was ERD present?]
% - [Summary of MRCP findings - was MRCP present?]
% - [Summary of artifact indicators - frontal activity?]
%
% Regarding what EEGNet learned:
% - Temporal filters: [target mu/beta? low frequency?]
% - Spatial filters: [focus on sensorimotor? frontal?]
%
% Regarding what drives classification:
% - Attribution analysis showed: [sensorimotor? frontal? other?]
% - Pre-cue vs post-cue: [which period dominates?]
%
% The dominant explanation:
% - [State the conclusion - which hypothesis best explains findings?]
% - [Note any significant secondary contributions]


% ============================================================
% 7.2 ANSWERS TO RESEARCH QUESTIONS
% ============================================================
\section{Answers to Research Questions}
\label{sec:concl_rqs}

% TARGET: 1-1.5 pages
%
% CONTENT:
%
% RQ1: What neural and artifactual signatures are present in the 
%      Apple Catcher dataset?
%
% Answer:
% - [Specific answer based on findings]
% - [Reference key results: ERDS analysis, MRCP, frontal activity]
%
% ----------------------------------------
%
% RQ2: Which features does EEGNet learn to extract?
%
% Answer:
% - [Temporal filters: which frequencies?]
% - [Spatial filters: which regions?]
% - [Reference key results: filter analysis]
%
% ----------------------------------------
%
% RQ3: Which features drive classification decisions, and are they 
%      neurophysiologically plausible?
%
% Answer:
% - [What attributions showed]
% - [Whether it aligns with neurophysiology]
% - [Whether it's plausible or concerning]


% ============================================================
% 7.3 CONTRIBUTIONS
% ============================================================
\section{Contributions}
\label{sec:concl_contributions}

% TARGET: 0.5-1 page
%
% CONTENT:
%
% This thesis makes the following contributions:
%
% 1. Systematic interpretability analysis of cross-subject EEGNet
%    - First such analysis for Apple Catcher paradigm
%    - Methodology applicable to other visual MI-BCI
%
% 2. Forensic investigation framework
%    - Competing hypotheses with testable predictions
%    - Evidence scoring methodology
%    - Template for validating BCI models
%
% 3. Evidence regarding Apple Catcher classification
%    - [Specific finding: what drives the 78% accuracy]
%    - [Actionable insight for deployment]
%
% 4. Insights for visual MI-BCI design
%    - [Lessons about visual cue confounds]
%    - [Recommendations for paradigm design]


% ============================================================
% 7.4 RECOMMENDATIONS
% ============================================================
\section{Recommendations}
\label{sec:concl_recommendations}

% TARGET: 0.5-1 page
%
% CONTENT:
%
% For Apple Catcher development:
% - [Based on findings: proceed? modify? reconsider?]
% - [Specific actionable steps]
%
% For visual MI-BCI researchers:
% - Always validate models with interpretability analysis
% - Check for artifact confounds, especially with visual cues
% - Don't assume accuracy equals validity
%
% For BCI deep learning practitioners:
% - Filter visualization and attribution should be standard
% - Compare learned features to expected neurophysiology
% - Report interpretability alongside accuracy


% ============================================================
% 7.5 FUTURE WORK
% ============================================================
\section{Future Work}
\label{sec:concl_future}

% TARGET: 0.5-1 page
%
% CONTENT:
%
% Immediate next steps:
% 1. [Most important follow-up experiment]
% 2. [Second priority]
%
% Longer-term directions:
% - [Broader research direction this work opens]
%
% Note: More detailed future directions in Discussion chapter


% ============================================================
% 7.6 CONCLUDING REMARKS
% ============================================================
\section{Concluding Remarks}
\label{sec:concl_final}

% TARGET: 0.5 page
%
% CONTENT:
%
% Big picture reflection:
% - Deep learning enables powerful BCI classification
% - But power without understanding is dangerous
% - Interpretability analysis is essential, not optional
%
% This thesis demonstrates:
% - How to systematically investigate what a model learns
% - That high accuracy requires validation, not celebration
% - [Final statement about findings and their significance]
